% Generated draft: requires mathematical and source review.
\documentclass[11pt]{article}
\usepackage[margin=25mm]{geometry}
\usepackage{fontspec}
\setmainfont{Latin Modern Roman}
\usepackage{amsmath,amssymb,mathtools,unicode-math}
\setmathfont{Latin Modern Math}
\usepackage{enumitem,needspace,xurl,hyperref}
\setlength{\parindent}{0pt}
\setlength{\parskip}{5pt}
\setlength{\emergencystretch}{3em}
\newcommand{\N}{\mathbb{N}}
\newcommand{\Z}{\mathbb{Z}}
\newcommand{\Q}{\mathbb{Q}}
\newcommand{\R}{\mathbb{R}}
\newcommand{\C}{\mathbb{C}}
\newcommand{\F}{\mathbb{F}}
\newcommand{\A}{\mathbb{A}}
\newcommand{\PP}{\mathbb{P}}
\newcommand{\E}{\mathbb{E}}
\newcommand{\eps}{\varepsilon}
\newcommand{\dd}{\,\mathrm{d}}
\begin{document}
\section*{UnsolvedMath ID: 30005524}
\textbf{Source problem code:} catalogue entry
\subsection{Definitions and mathematical statement}
Bone mineral density (BMD) is the quantity whose measurement is at issue. DXA denotes dual-energy X-ray absorptiometry, and CT denotes computed tomography. The supplied record does not specify a mathematical definition of the measured density, its units, or whether an areal or volumetric quantity is intended.

The source question is: \emph{Are bone-mineral-density measurements obtained by DXA or CT reliable?}

This is an unresolved formulation of a measurement-reliability question rather than a fully specified mathematical proposition. The record supplies no measurement operator, reference quantity, calibration procedure, error model, acquisition protocol, anatomical region, patient population, or quantitative acceptance criterion. In particular, it does not define whether reliability means accuracy relative to a reference, repeatability, reproducibility, diagnostic performance, or agreement between modalities. It also does not specify whether DXA and CT are to be evaluated separately or compared for interchangeability. Thus no definite inequality, domain, boundary conditions, or quantified reliability assertion can be extracted without additional hypotheses.

The supplied dated assessment distinguishes quantitative CT from CT in general: it describes DXA as the clinical reference with established thresholds and says that quantitative CT can be reliable after calibration and validation, but is not generally interchangeable with DXA. These are supplied curation claims, not mathematical conclusions established by this draft. The record does not provide the thresholds, calibration assumptions, validation criteria, or scope needed to incorporate those claims into a precise mathematical statement.
\subsection{Short English statement}
The question asks whether DXA or CT measurements of bone mineral density are reliable. A precise mathematical statement requires a specified measured quantity, measurement setting, reference standard, and definition of reliability. The supplied assessment is dated 2026-08-22 and labels the problem partially solved; that assessment has not been independently verified.
\subsection{Sources}
\begin{itemize}
\item Local UnsolvedMath registry: \texttt{data/unsolvedmath/problems.json}, with \texttt{catalog/problems.json} as the supplementary registry; permanent UnsolvedMath ID \(30005524\), title \emph{Reliability of Bone Mineral Density Imaging}, source problem code \texttt{catalogue entry}. The supplied record belongs to the Oberwolfach Reports Open Problems collection and gives proposed year 2023; no individual proposer is supplied.
\item Supplied report attribution: \emph{Tomographic Inverse Problems: Mathematical Challenges and Novel Applications} (2023), Oberwolfach Reports, pp.\ 1105--1194, DOI 10.4171/owr/2023/21. \url{https://doi.org/10.4171/owr/2023/21}. The background describes Oberwolfach Reports as documenting mathematical workshops and their open-problem discussions. No specific page or speaker for this question is identified.
\item Supplied literature reference: ISCD official positions. \url{https://iscd.org/official-positions-2023/}.
\item Supplied literature reference: CT and DXA validation study. \url{https://doi.org/10.1038/s41598-024-79985-w}.
\item Literature assessment date: \texttt{2026-08-22}. This is supplied dated curation, not an independently verified assessment. Its status label is \texttt{partially\_solved}. It states that DXA remains the clinical reference with established thresholds, while quantitative CT can be reliable after calibration and validation but is not generally interchangeable with DXA. The background expressly calls this dated literature triage derived from the supplied OWR curation and requests independent verification of time-sensitive status and cited claims.
\item The registry labels the statement \texttt{corrected\_verified} and reports that the official Oberwolfach report was checked directly and that the clean statement preserves the complete mathematical target while restoring definitions and removing extraction artifacts or source-dependent wording. These are supplied provenance assertions; this draft does not independently verify them.
\end{itemize}
\end{document}
